% AA-18 | Algebra abstrakcyjna | Iteracja IX
\chapter{Iloczyn tensorowy modułów i jego własności}
\label{aa:ch:module-tensor-product}

\section{Dlaczego potrzebujemy nowej konstrukcji?}
Rozważmy dwie grupy abelowe $M,N$. Iloczyn prosty $M\times N$
przechowuje po prostu parę ich elementów. Iloczyn tensorowy jest
czymś innym: pozwala zamienić \emph{dwuliniowe} odwzorowania
z $M\times N$ na homomorfizmy z jednej grupy. W szczególności
iloczyn tensorowy dwóch niezerowych modułów może być zerowy.

Ważne jest rozróżnienie trzech konstrukcji. W iloczynie prostym
występują uporządkowane pary $(m,n)$. W sumie prostej występują
rodziny elementów o skończonym nośniku. W iloczynie tensorowym
występują klasy formalnych symboli $m\otimes n$, podlegające
konkretnym relacjom addytywności i zgodności ze skalarami.
Nie utożsamiamy żadnej z tych konstrukcji z pozostałymi.

Jeśli pierścień $R$ nie jest przemienny, strony działania
skalarów mają znaczenie. Do podstawowej konstrukcji nad $R$
potrzebujemy \emph{prawego} $R$-modułu $M$ i \emph{lewego}
$R$-modułu $N$. Otrzymany iloczyn tensorowy jest w ogólności
grupą abelową; bez dodatkowych struktur nie musi mieć
kanonicznego działania pierścienia $R$.

\section{Odwzorowania zbilansowane i konstrukcja ilorazowa}
\begin{definicja}[Odwzorowanie addytywne i zbilansowane]
\label{aa:def:tensor-balanced}
Niech $R$ będzie pierścieniem z jedynką, niekoniecznie
przemiennym, $M$ prawym $R$-modułem, $N$ lewym $R$-modułem,
a $A$ grupą abelową. Odwzorowanie $b:M\times N\to A$
nazywamy \emph{addytywnym w każdej zmiennej i zbilansowanym},
gdy dla wszystkich $m,m'\in M$, $n,n'\in N$ i $r\in R$ zachodzą
\[
\begin{split}
b(m+m',n)&=b(m,n)+b(m',n),\\
b(m,n+n')&=b(m,n)+b(m,n'),\\
b(mr,n)&=b(m,rn).
\end{split}
\]
Trzeci warunek pozwala przenieść skalar działający
z prawej strony na $M$ do działania z lewej strony na $N$.
Jest konieczny do zdefiniowania iloczynu tensorowego nad $R$.
\end{definicja}

\begin{twierdzenie}[Istnienie iloczynu tensorowego i własność uniwersalna]
\label{aa:tw:tensor-universal}
Niech $R$ będzie pierścieniem z jedynką, niekoniecznie
przemiennym, $M$ prawym $R$-modułem, a $N$ lewym
$R$-modułem. Wtedy istnieje grupa abelowa $M\otimes_R N$ oraz odwzorowanie
$\tau:M\times N\to M\otimes_R N$ addytywne w każdej
zmiennej i zbilansowane, takie, że dla każdej grupy abelowej
$A$ i każdego odwzorowania $b:M\times N\to A$ o tych
własnościach istnieje dokładnie jeden homomorfizm grup
$\bar b:M\otimes_R N\to A$ z $\bar b\circ\tau=b$.
Ponadto para $(M\otimes_RN,\tau)$ jest wyznaczona
jednoznacznie z dokładnością do jedynego izomorfizmu
zgodnego z $\tau$.
\end{twierdzenie}
\begin{proof}
Niech $F$ będzie \emph{wolną grupą abelową} na zbiorze
symboli $[m,n]$, gdzie $(m,n)\in M\times N$.
Oznacza to, że elementy $F$ są formalnymi skończonymi
sumami całkowitych wielokrotności tych symboli; różne
symbole są niezależnymi generatorami.

Niech $L$ będzie podgrupą $F$ generowaną przez elementy
\[
\begin{split}
 &[m+m',n]-[m,n]-[m',n],\\
 &[m,n+n']-[m,n]-[m,n'],\\
 &[mr,n]-[m,rn]
\end{split}
\]
dla wszystkich dopuszczalnych $m,m',n,n',r$.
Definiujemy $M\otimes_RN=F/L$ i
$m\otimes n=[m,n]+L$. Wymienione relacje zapewniają,
że $(m,n)\mapsto m\otimes n$ spełnia dokładnie
trzy wymagane warunki. Z addytywności wynika też
$0\otimes n=m\otimes0=0$ oraz
$(-m)\otimes n=-(m\otimes n)$.

Niech teraz $b:M\times N\to A$ będzie jak w tezie.
Ponieważ $F$ jest wolną grupą abelową na zbiorze par,
istnieje dokładnie jeden homomorfizm $B:F\to A$
taki, że $B([m,n])=b(m,n)$. Trzy założenia o $b$
mówią, że $B$ zeruje wszystkie generatory $L$.
Zatem $L\subseteq\ker B$ i wzór
$\bar b(x+L)=B(x)$ określa homomorfizm na ilorazie.
Z definicji $\bar b(m\otimes n)=b(m,n)$.
Każdy element $F/L$ jest skończoną sumą całkowitych
wielokrotności tensorów prostych; stąd homomorfizm
o wskazanych wartościach jest jednoznaczny.

Jeśli $(T,t)$ jest inną parą o tej samej własności,
uniwersalność daje homomorfizmy $T\to M\otimes_RN$
oraz $M\otimes_RN\to T$, zachowujące obrazy wszystkich
par $(m,n)$. Ich złożenia zachowują te same generatory,
więc są identycznościami. Homomorfizmy są odwrotne
i dają wymagany jednoznaczny izomorfizm.
\end{proof}

\begin{uwaga}[Oznaczenie tensoru prostego]
Element $m\otimes n$ nazywamy \emph{tensorem prostym}.
Nie każdy element $M\otimes_RN$ musi być tensorem prostym:
na ogół jest sumą takich elementów. Na przykład dla dowolnego ciała $K$, przestrzeni
$V=W=K^2$ i baz $e_1,e_2$ oraz $f_1,f_2$ tensor
$e_1\otimes f_1+e_2\otimes f_2$ nie jest prosty.
Gdyby był równy $(ae_1+be_2)\otimes(cf_1+df_2)$,
to porównanie współczynników przy czterech tensorach
bazowych dawałoby $ac=bd=1$ oraz $ad=bc=0$.
Pierwsze dwie równości wymuszają $a,b,c,d\ne0$,
co przeczy pozostałym dwóm. Wykorzystaliśmy tu
niezależność tensorów bazowych, udowodnioną
w Stwierdzeniu~\ref{aa:st:tensor-vspace-dimension}.
Równość $m\otimes n=m'\otimes n'$ nie oznacza, że
$m=m'$ i $n=n'$. W szczególności z $m\otimes n=0$
nie wynika, że któryś z czynników jest zerowy;
porównajmy Przykład~\ref{aa:prz:tensor-cyclic}.
\end{uwaga}

\section{Przemienny pierścień współczynników}
\begin{twierdzenie}[Iloczyn tensorowy jako moduł]
\label{aa:tw:tensor-commutative-module}
Niech $R$ będzie pierścieniem przemiennym z jedynką,
a $M,N$ lewymi $R$-modułami. Traktując $M$ jako prawy
moduł przez $mr=rm$, otrzymujemy $R$-moduł
$M\otimes_R N$ z działaniem
$r(m\otimes n)=(rm)\otimes n=m\otimes(rn)$.
Dla każdego $R$-modułu $P$ odwzorowania $R$-dwuliniowe
$M\times N\to P$ odpowiadają bijektywnie homomorfizmom
$R$-modułów $M\otimes_RN\to P$.
\end{twierdzenie}
\begin{proof}
Dla $r\in R$ na wolnej grupie abelowej $F$ z poprzedniego
dowodu określmy $T_r([m,n])=[rm,n]$ i przedłużmy
addytywnie. Sprawdzimy, że $T_r(L)\subseteq L$.
Dla pierwszego rodzaju relacji addytywnych otrzymujemy
$[r(m+m'),n]-[rm,n]-[rm',n]$; z równości
$r(m+m')=rm+rm'$ wynika, że jest to generator
pierwszego rodzaju. Dla drugiego rodzaju
otrzymujemy bezpośrednio relację
$[rm,n+n']-[rm,n]-[rm,n']$. Dla relacji zbilansowania
obraz elementu $[ms,n]-[m,sn]$ wynosi
$[r(ms),n]-[rm,sn]=[(rm)s,n]-[rm,sn]$,
ponieważ $R$ jest przemienny. Jest to generator $L$.
Zatem $T_r$ przechodzi na grupę ilorazową.
Wzory $T_{r+s}=T_r+T_s$, $T_{rs}=T_rT_s$
i $T_1=\operatorname{id}$ sprawdzamy na generatorach
$[m,n]$, używając odpowiednio addytywności działania,
$(rs)m=r(sm)$ oraz $1_Rm=m$. Przechodzą one na
iloraz i dają wszystkie aksjomaty $R$-modułu.
Równość $(rm)\otimes n=m\otimes(rn)$ jest
jedną z relacji zbilansowania.

Jeśli $b:M\times N\to P$ jest $R$-dwuliniowe,
jest addytywne w każdej zmiennej i zbilansowane,
gdyż $b(rm,n)=rb(m,n)=b(m,rn)$.
Własność uniwersalna daje homomorfizm grup
$\bar b:M\otimes_RN\to P$. Na tensorach prostych
$\bar b(r(m\otimes n))=b(rm,n)=rb(m,n)$,
więc $\bar b$ jest $R$-liniowy. Odwrotnie, złożenie
$R$-liniowego homomorfizmu z kanonicznym odwzorowaniem
$(m,n)\mapsto m\otimes n$ jest $R$-dwuliniowe.
Obie konstrukcje są wzajemnie odwrotne.
\end{proof}

\begin{przyklad}[Iloczyn tensorowy grup abelowych]
\label{aa:prz:tensor-abelian-example}
Każda grupa abelowa jest $\mathbb Z$-modułem.
Zatem dla grup abelowych $A,B$ zapis
$A\otimes_{\mathbb Z}B$ oznacza ich iloczyn tensorowy
jako modułów nad $\mathbb Z$. Jest to dokładnie
ta sama konstrukcja, którą wprowadziliśmy w
Rozdziale~\ref{aa:ch:grupy-abelowe}, teraz
umieszczona w ogólnej teorii modułów.
\end{przyklad}

\section{Podstawowe izomorfizmy}
\begin{twierdzenie}[Element neutralny dla iloczynu tensorowego]
\label{aa:tw:tensor-unit}
Niech $R$ będzie pierścieniem z jedynką, a $N$
lewym $R$-modułem. Odwzorowanie
$R\otimes_RN\to N$, $r\otimes n\mapsto rn$,
jest izomorfizmem grup abelowych. Jeśli $R$ jest
przemienny, jest też izomorfizmem $R$-modułów.
Analogicznie dla prawego $R$-modułu $M$ mamy
$M\otimes_RR\cong M$.
\end{twierdzenie}
\begin{proof}
Odwzorowanie $(r,n)\mapsto rn$ jest addytywne
w każdej zmiennej i zbilansowane: $(rs)n=r(sn)$.
Uniwersalność daje homomorfizm $F:R\otimes_RN\to N$.
Zdefiniujmy $G:N\to R\otimes_RN$ przez $G(n)=1\otimes n$.
Jest to homomorfizm grup. Wtedy
$F(G(n))=1n=n$, natomiast
$G(F(r\otimes n))=1\otimes(rn)=r\otimes n$.
Ostatnia równość wynika z relacji zbilansowania.
Ponieważ tensory proste generują całą grupę,
$G\circ F=\operatorname{id}$. Dla przemiennego $R$
oba homomorfizmy zachowują działanie skalarów.
Dla prawego $R$-modułu $M$ druga teza wynika wprost
z odwzorowań $m\otimes r\mapsto mr$ oraz
$m\mapsto m\otimes1_R$. Ich złożenia są identycznościami,
ponieważ $(mr)\otimes1_R=m\otimes(r1_R)=m\otimes r$.
\end{proof}

\begin{twierdzenie}[Suma prosta i iloczyn tensorowy]
\label{aa:tw:tensor-direct-sum}
Niech $R$ będzie przemiennym pierścieniem z jedynką,
a $M,N,P$ jego modułami. Wówczas istnieje kanoniczny izomorfizm
\[
 (M\oplus N)\otimes_R P\cong
 (M\otimes_RP)\oplus(N\otimes_RP).
\]
Taki sam wzór zachodzi dla dowolnej rodziny modułów
z sumą prostą po lewej i sumą prostą ich iloczynów
tensorowych po prawej.
\end{twierdzenie}
\begin{proof}
Dwuliniowe odwzorowanie
$((m,n),p)\mapsto(m\otimes p,n\otimes p)$
indukuje homomorfizm $F$ z lewej strony na prawą.
W odwrotnym kierunku definiujemy homomorfizmy
$M\otimes_RP\to(M\oplus N)\otimes_RP$ przez
$m\otimes p\mapsto(m,0)\otimes p$ oraz
$N\otimes_RP\to(M\oplus N)\otimes_RP$ przez
$n\otimes p\mapsto(0,n)\otimes p$.
Są one poprawne dzięki własności uniwersalnej.
Ich suma daje homomorfizm $G$ z prawej strony na lewą.
Na tensorach prostych mamy
$GF((m,n)\otimes p)=(m,0)\otimes p+(0,n)\otimes p
=(m,n)\otimes p$.
Z drugiej strony $FG$ zachowuje osobno oba składniki,
więc jest identycznością. Dla dowolnej rodziny $(M_i)_{i\in I}$ definiujemy
$F((m_i)_{i\in I}\otimes p)=(m_i\otimes p)_{i\in I}$.
Wyrażenie po prawej ma skończony nośnik, gdyż
$(m_i)$ należy do sumy prostej. Jest dwuliniowe,
więc daje homomorfizm z iloczynu tensorowego.
W przeciwną stronę każda kanoniczna inkluzja
$M_i\hookrightarrow\bigoplus_{j\in I}M_j$ daje
homomorfizm $M_i\otimes_RP\to
(\bigoplus_{j\in I}M_j)\otimes_RP$.
Dla rodziny tensorów o skończonym nośniku sumujemy
ich obrazy po tych inkluzjach; otrzymujemy $G$.
Na tensorach prostych obie kompozycje są
identycznościami, a na prawej sumie prostej
sprawdzamy to oddzielnie na każdym składniku.
Zatem $F$ i $G$ są wzajemnie odwrotne.
\end{proof}

\begin{twierdzenie}[Przemienność i łączność iloczynu tensorowego]
\label{aa:tw:tensor-sym-assoc}
Niech $R$ będzie przemiennym pierścieniem z jedynką,
a $M,N,P$ jego modułami. Istnieją kanoniczne izomorfizmy
\[
 M\otimes_RN\cong N\otimes_RM,
 \qquad (M\otimes_RN)\otimes_RP\cong
 M\otimes_R(N\otimes_RP),
\]
wyznaczone odpowiednio przez $m\otimes n\mapsto n\otimes m$
oraz $(m\otimes n)\otimes p\mapsto m\otimes(n\otimes p)$.
\end{twierdzenie}
\begin{proof}
Odwzorowanie $(m,n)\mapsto n\otimes m$ jest
$R$-dwuliniowe, więc daje homomorfizm dla pierwszego
izomorfizmu. To samo odwzorowanie w odwrotnej kolejności
jest jego odwrotnością na generatorach.

Przy drugim izomorfizmie sprawdzenie poprawności wymaga
dwóch kolejnych zastosowań własności uniwersalnej.
Dla ustalonego $p$ odwzorowanie $(m,n)\mapsto
m\otimes(n\otimes p)$ jest $R$-dwuliniowe, więc
daje homomorfizm $\alpha_p:M\otimes_RN\to
M\otimes_R(N\otimes_RP)$. Dla ustalonego $x$ odwzorowanie
$p\mapsto\alpha_p(x)$ jest liniowe: dla
$x=m\otimes n$ wynika to z liniowości
$p\mapsto n\otimes p$, a ogólny $x$ jest sumą
skończenie wielu tensorów prostych. Podobnie
$\alpha_{rp}(m\otimes n)=r\alpha_p(m\otimes n)$
i $\alpha_p(rx)=r\alpha_p(x)$, zatem
$(x,p)\mapsto\alpha_p(x)$ jest $R$-dwuliniowe.
Przez kolejne zastosowanie własności uniwersalnej
otrzymujemy homomorfizm
$(M\otimes_RN)\otimes_RP\to M\otimes_R(N\otimes_RP)$.
W odwrotną stronę konstruujemy analogicznie
$m\otimes(n\otimes p)\mapsto(m\otimes n)\otimes p$.
Na elementach $(m\otimes n)\otimes p$ i
$m\otimes(n\otimes p)$ złożenia tych homomorfizmów
są identycznościami; ponieważ generują one odpowiednie
moduły, homomorfizmy są wzajemnie odwrotne.
\end{proof}

\section{Iloczyn tensorowy z modułem ilorazowym}
\begin{twierdzenie}[Iloczyn tensorowy z ilorazem pierścienia]
\label{aa:tw:tensor-quotient}
Niech $R$ będzie przemiennym pierścieniem z jedynką,
$I\subseteq R$ ideałem, a $N$ lewym $R$-modułem.
Przez $IN$ oznaczamy podmoduł generowany przez wszystkie
$in$, gdzie $i\in I$ i $n\in N$. Wtedy istnieje naturalny
izomorfizm $R$-modułów
\[
 (R/I)\otimes_R N\cong N/IN,
 \qquad (r+I)\otimes n\longmapsto rn+IN.
\]
\end{twierdzenie}
\begin{proof}
Definiujemy odwzorowanie $b:(R/I)\times N\to N/IN$
przez $b(r+I,n)=rn+IN$. Jeśli $r'=r+i$ dla $i\in I$,
to $(r'-r)n=in\in IN$, więc wzór nie zależy od
reprezentanta $r+I$. Zachowuje dodawanie w obu
zmiennych i spełnia $b((r+I)s,n)=b(r+I,sn)$.
Dlatego daje homomorfizm $F$ z iloczynu tensorowego
na $N/IN$.

W drugą stronę definiujemy $G:N/IN\to (R/I)\otimes_RN$
przez $G(n+IN)=(1+I)\otimes n$.
Jeżeli $n'=n+\sum_j i_j n_j$ z $i_j\in I$, to
$G(n')-G(n)=\sum_j(1+I)\otimes i_jn_j
=\sum_j(i_j+I)\otimes n_j=0$.
Zatem $G$ jest dobrze określony i $R$-liniowy.
Dla $n\in N$ mamy $FG(n+IN)=n+IN$.
Dla generatora $(r+I)\otimes n$ z drugiej strony
$GF((r+I)\otimes n)=(1+I)\otimes rn
=(r+I)\otimes n$.
Obie kompozycje są więc identycznościami.
\end{proof}

\begin{wniosek}[Dwa ilorazy pierścienia]
\label{aa:wn:tensor-two-quotients}
Niech $R$ będzie przemiennym pierścieniem z jedynką,
a $I,J\subseteq R$ jego ideałami. Wówczas istnieje
naturalny izomorfizm $R$-modułów
$(R/I)\otimes_R(R/J)\cong R/(I+J)$.
\end{wniosek}
\begin{proof}
Z Twierdzenia~\ref{aa:tw:tensor-quotient} dostajemy
$(R/I)\otimes_R(R/J)\cong(R/J)/I(R/J)$.
Obrazem podmodułu $I(R/J)$ są dokładnie klasy
elementów ideału $I+J$ modulo $J$.
Trzecie twierdzenie o izomorfizmie modułów daje
$(R/J)/((I+J)/J)\cong R/(I+J)$.
\end{proof}

\begin{przyklad}[Iloczyn tensorowy grup cyklicznych]
\label{aa:prz:tensor-cyclic}
Dla liczb całkowitych $m,n\ge1$ oraz $d=\gcd(m,n)$
zachodzi
\[
 (\mathbb Z/m\mathbb Z)\otimes_{\mathbb Z}
 (\mathbb Z/n\mathbb Z)\cong \mathbb Z/d\mathbb Z.
\]
Rzeczywiście $(m)+(n)=(d)$, więc wystarczy
zastosować Wniosek~\ref{aa:wn:tensor-two-quotients}.
Dla $m=12$, $n=18$ otrzymujemy grupę cykliczną
rzędu sześć. Natomiast dla $m=2$, $n=3$ iloczyn
tensorowy jest zerowy. Jest to wyraźnie różna sytuacja
od iloczynu prostego $C_2\times C_3\cong C_6$.
\end{przyklad}

\begin{przyklad}[Iloczyn tensorowy z ciałem liczb wymiernych]
\label{aa:prz:tensor-rational-torsion}
Dla $n\ge1$ mamy $\mathbb Q\otimes_{\mathbb Z}
(\mathbb Z/n\mathbb Z)=0$. Istotnie elementy
$\mathbb Q$ można dzielić przez $n$, więc
$n\mathbb Q=\mathbb Q$. Z Twierdzenia
\ref{aa:tw:tensor-quotient}, po zamianie czynników,
otrzymujemy $(\mathbb Z/n\mathbb Z)\otimes_{\mathbb Z}
\mathbb Q\cong\mathbb Q/n\mathbb Q=0$.
Nie oznacza to, że $\mathbb Q$ lub $\mathbb Z/n\mathbb Z$
są modułami zerowymi; zerowy jest ich iloczyn tensorowy.
\end{przyklad}

\section{Iloczyn tensorowy i krótkie ciągi dokładne}
\begin{twierdzenie}[Prawostronna dokładność funktora iloczynu tensorowego]
\label{aa:tw:tensor-right-exact}
Niech $R$ będzie przemiennym pierścieniem z jedynką,
a $N$ jego modułem. Jeśli ciąg homomorfizmów
$A\xrightarrow{f}B\xrightarrow{g}C\to0$ jest dokładny,
to dokładny jest również ciąg
\[
 A\otimes_RN\xrightarrow{f\otimes\operatorname{id}_N}
 B\otimes_RN\xrightarrow{g\otimes\operatorname{id}_N}
 C\otimes_RN\longrightarrow0.
\]
Dla dowolnego homomorfizmu $f$ odwzorowanie
$f\otimes\operatorname{id}_N$ jest określone przez
$a\otimes n\mapsto f(a)\otimes n$.
\end{twierdzenie}
\begin{proof}
Najpierw poprawność ostatniego wzoru wynika z własności
uniwersalnej, ponieważ $(a,n)\mapsto f(a)\otimes n$
jest dwuliniowe. Z dokładności pierwotnego ciągu
mamy $g$ surjektywne i $\ker g=\operatorname{im}f$.
Każdy tensor prosty $c\otimes n\in C\otimes_RN$
ma przedobraz $b\otimes n$ z $g(b)=c$; ponieważ
tensory proste generują, odwzorowanie
$g\otimes\operatorname{id}_N$ jest surjektywne.

Oznaczmy $I=\operatorname{im}f\subseteq B$.
Z pierwszego twierdzenia o izomorfizmie modułów
mamy $C\cong B/I$. Rozważmy teraz iloraz
$D=(B\otimes_RN)/H$, gdzie $H$ jest podmodułem
generowanym przez wszystkie $i\otimes n$ z
$i\in I$, $n\in N$. Odwzorowanie
$(b+I,n)\mapsto b\otimes n+H$
jest dobrze określone: zmiana $b$ o $i\in I$
zmienia jego wartość o $i\otimes n\in H$.
Jest dwuliniowe, więc otrzymujemy
$(B/I)\otimes_RN\to D$.
Z drugiej strony odwzorowanie
$b\otimes n\mapsto(b+I)\otimes n$ zeruje $H$
i daje homomorfizm $D\to(B/I)\otimes_RN$.
Oba złożenia zachowują tensory proste, a zatem są
identycznościami. Tak więc
$D\cong(B/I)\otimes_RN\cong C\otimes_RN$.
Oznacza to, że jądro $g\otimes\operatorname{id}_N$
jest dokładnie $H$. Ale $H$ jest obrazem
$f\otimes\operatorname{id}_N$: obraz tego homomorfizmu
jest podmodułem, zawiera każdy generator $i\otimes n$
modułu $H$, gdyż $i=f(a)$ dla pewnego $a\in A$,
a jednocześnie obrazy wszystkich tensorów prostych
$a\otimes n$ należą do $H$. Ponieważ tensory proste
generują $A\otimes_RN$, oba podmoduły są równe.
Dostaliśmy żądaną dokładność.
\end{proof}

\begin{przyklad}[Dlaczego iniekcja może zniknąć]
\label{aa:prz:tensor-not-left-exact}
Rozważmy krótki ciąg dokładny $\mathbb Z$-modułów
\[
0\longrightarrow\mathbb Z\xrightarrow{\cdot2}\mathbb Z
\longrightarrow\mathbb Z/2\mathbb Z\longrightarrow0.
\]
Po wzięciu iloczynu tensorowego z $\mathbb Z/2\mathbb Z$
i użyciu $\mathbb Z\otimes_{\mathbb Z}\mathbb Z/2\mathbb Z
\cong\mathbb Z/2\mathbb Z$ otrzymujemy odwzorowanie
$\mathbb Z/2\mathbb Z\xrightarrow{\cdot2}
\mathbb Z/2\mathbb Z$. Jest ono zerowe, ponieważ
$(2a)\otimes\bar b=a\otimes(2\bar b)=0$ dla
każdych $a\in\mathbb Z$, $\bar b\in\mathbb Z/2\mathbb Z$.
Nie jest ono iniekcją. W ogólności \emph{iloczyn tensorowy
nie zachowuje iniektywności}, choć zachowuje surjektywność.
To odróżnia go również od iloczynu prostego modułów.
\end{przyklad}

\section{Moduły płaskie i moduły rzutowe}
\begin{definicja}[Moduł płaski]
\label{aa:def:flat-module}
Niech $R$ będzie przemiennym pierścieniem z jedynką.
$R$-moduł $N$ nazywamy \emph{płaskim}, jeżeli dla
każdego krótkiego ciągu dokładnego
$0\to A\to B\to C\to0$ ciąg uzyskany przez
wzięcie iloczynu tensorowego z $N$,
\[
0\longrightarrow A\otimes_RN\longrightarrow B\otimes_RN
\longrightarrow C\otimes_RN\longrightarrow0,
\]
jest także dokładny.
\end{definicja}

\begin{twierdzenie}[Moduły wolne i rzutowe są płaskie]
\label{aa:tw:projective-flat}
Niech $R$ będzie przemiennym pierścieniem z jedynką.
Każdy wolny $R$-moduł jest płaski; każdy składnik sumy
prostej modułu płaskiego jest płaski. W szczególności
każdy moduł rzutowy jest płaski.
\end{twierdzenie}
\begin{proof}
Jeśli $F=R^{(S)}$ jest modułem wolnym na zbiorze $S$,
rozdzielność iloczynu tensorowego względem sum prostych
z Twierdzenia~\ref{aa:tw:tensor-direct-sum} i izomorfizm
$M\otimes_RR\cong M$ pokazują, że
$M\otimes_RF\cong M^{(S)}$.
Suma prosta kopii krótkiego ciągu dokładnego jest
dokładna współrzędnie: element sumy prostej ma skończony
nośnik, więc badanie jąder i obrazów odbywa się w skończonej
liczbie współrzędnych. Zatem $F$ jest płaski.

Załóżmy, że $N=P\oplus Q$ jest płaski. Dla każdej
iniekcji $f:A\to B$ homomorfizm
$f\otimes\operatorname{id}_N$ jest iniektywny.
Pod identyfikacjami
$A\otimes_RN\cong(A\otimes_RP)\oplus(A\otimes_RQ)$
i odpowiednio dla $B$ jest on sumą prostą
$f\otimes\operatorname{id}_P$ oraz
$f\otimes\operatorname{id}_Q$. Każda z tych map
musi być iniektywna, ponieważ ich suma prosta jest
iniektywna. W połączeniu z prawostronną dokładnością
dowodzi to płaskości $P$ oraz $Q$.
W Rozdziale~\ref{aa:ch:modules-introduction} udowodniliśmy,
że moduł rzutowy jest składnikiem sumy prostej modułu
wolnego; dlatego każdy moduł rzutowy jest płaski.
\end{proof}

\begin{przyklad}[Moduł niepłaski]
Moduł $\mathbb Z/2\mathbb Z$ nie jest płaski nad
$\mathbb Z$, ponieważ przykład
\ref{aa:prz:tensor-not-left-exact} pokazuje,
że iloczyn tensorowy z nim nie zachowuje iniekcji.
\end{przyklad}

\begin{twierdzenie}[Ciało ułamków dziedziny całkowitości jest modułem płaskim]
\label{aa:tw:frac-flat}
Niech $R$ będzie dziedziną całkowitości i $K$
jej ciałem ułamków. Wtedy $K$ jest płaskim
$R$-modułem.
\end{twierdzenie}
\begin{proof}
Użyjemy jawnego opisu lokalizacji. Dla dowolnego
$R$-modułu $M$ niech $S=R\setminus\{0\}$.
Na parach $(m,s)\in M\times S$ przyjmujemy
$(m,s)\sim(m',s')$ wtedy i tylko wtedy, gdy
istnieje $t\in S$ takie, że $t(s'm-sm')=0$.
Jest to relacja równoważności: zwrotność i symetria
są natychmiastowe. Dla przechodniości załóżmy, że
$t(s'm-sm')=0$ i $u(s''m'-s'm'')=0$.
Po pomnożeniu pierwszej równości przez $us''$,
a drugiej przez $ts$, i dodaniu otrzymujemy
$tu s'(s''m-sm'')=0$; element $tu s'$ należy do $S$.

Na zbiorze klas $S^{-1}M$ określamy działania
$[m,s]+[m',s']=[s'm+sm',ss']$ oraz
$r[m,s]=[rm,s]$. Są one niezależne od reprezentantów.
Istotnie, z $t(s_1m-sm_1)=0$ i
$u(s'_1m'-s'm'_1)=0$ wynika
\[
\begin{aligned}
&s_1s'_1(s'm+sm')-ss'(s'_1m_1+s_1m'_1)\\
&\quad=s'_1s'(s_1m-sm_1)+s_1s(s'_1m'-s'm'_1).
\end{aligned}
\]
Po pomnożeniu prawej strony przez $tu\in S$
otrzymujemy zero. Jest to dokładnie warunek równoważności
sum dwóch par. Analogicznie $t(s_1rm-srm_1)=0$
uzasadnia poprawność mnożenia klasy przez $r$.
Łączność dodawania wynika z tego, że suma trzech
klas, niezależnie od nawiasowania, ma przedstawiciela
$[s_2s_3m_1+s_1s_3m_2+s_1s_2m_3,s_1s_2s_3]$.
Przemienność dodawania wynika z przemienności $M$.
Elementem zerowym jest $[0,1]$, bo
$[m,s]+[0,1]=[m,s]$, a przeciwnym do $[m,s]$
jest $[-m,s]$, gdyż $[0,s^2]=[0,1]$.
Dla $r,q\in R$ i klas $x=[m,s]$, $y=[n,t]$
zachodzą $r(x+y)=rx+ry$ oraz
$(r+q)x=rx+qx$, co otrzymujemy przez sprowadzenie
do wspólnego mianownika. Ponadto
$r(qx)=[rqm,s]=(rq)x$ i $1_Rx=x$.
Wszystkie aksjomaty $R$-modułu są więc spełnione.
Dla homomorfizmu $f:M\to N$ wzór
$S^{-1}f([m,s])=[f(m),s]$ jest poprawny,
ponieważ $f(t(s'm-sm'))=t(s'f(m)-sf(m'))$.

Przypomnijmy, że $K=S^{-1}R$ jest ciałem ułamków.
Własność uniwersalna iloczynu tensorowego daje
homomorfizm $K\otimes_RM\to S^{-1}M$ zadany przez
$(r/s)\otimes m\mapsto[rm,s]$. Wzór jest poprawny
względem reprezentanta $r/s$: z $rs'=r's$
wynika $[rm,s]=[r'm,s']$. Addytywność względem
pierwszego argumentu sprawdzamy na wspólnym mianowniku:
para $(r/s+r'/s',m)$ jest wysyłana na
$[(rs'+r's)m,ss']$, a klasa ta jest równa
$[rm,s]+[r'm,s']$.
Addytywność względem $m$ wynika z
$[r(m+m'),s]=[rm,s]+[rm',s]$.
Zbilansowanie wynika z $[(rq)m,s]=[r(qm),s]$
dla każdego $q\in R$. Własność uniwersalna daje
zatem homomorfizm grup abelowych.
Odwrotność definiujemy przez
$[m,s]\mapsto(1/s)\otimes m$.
Jeśli $[m,s]=[m',s']$, to dla pewnego $t\in S$
mamy $t(s'm-sm')=0$. Wtedy
$(1/s)\otimes m=(1/(tss'))\otimes(t s'm)$ oraz
$(1/s')\otimes m'=(1/(tss'))\otimes(t sm')$,
więc wzór nie zależy od reprezentanta.
Jest addytywny, gdyż
$(1/(ss'))\otimes(s'm+sm')=
(1/s)\otimes m+(1/s')\otimes m'$, i zachowuje
mnożenie przez skalary. Oba złożenia zachowują
elementy postaci $[m,s]$ i $(r/s)\otimes m$,
a więc są identycznościami.
Sprawdźmy ponadto naturalność tego izomorfizmu.
Oznaczmy go przez $\Theta_M:K\otimes_RM\to S^{-1}M$.
Dla każdego homomorfizmu $f:M\to N$ i tensora prostego
$(r/s)\otimes m$ zachodzi
$\Theta_N((\operatorname{id}_K\otimes f)((r/s)\otimes m))
=[rf(m),s]=S^{-1}f(\Theta_M((r/s)\otimes m))$.
Ponieważ tensory proste generują $K\otimes_RM$,
otrzymujemy równość homomorfizmów
$\Theta_N\circ(\operatorname{id}_K\otimes f)
=(S^{-1}f)\circ\Theta_M$.

Zobaczmy, że lokalizacja zachowuje dokładność.
Jeżeli $0\to A\xrightarrow{f}B\xrightarrow{g}C\to0$
jest dokładny, to $S^{-1}g$ jest surjektywny:
każdy $[c,s]$ ma przedobraz $[b,s]$ dla pewnego
$b$ z $g(b)=c$. Jeśli $[b,s]$ przechodzi na zero,
to istnieje $t\in S$ takie, że $t g(b)=0$,
a więc $tb\in\ker g=\operatorname{im}f$.
Wybierzmy $a$ z $f(a)=tb$. Wtedy
$[b,s]=[tb,ts]=[f(a),ts]$, czyli klasa należy
do obrazu $S^{-1}f$. Wreszcie, jeśli
$[a,s]\in S^{-1}A$ przechodzi na zero pod
$S^{-1}f$, to $tf(a)=0$ dla pewnego $t\in S$.
Iniektywność $f$ daje $ta=0$, czyli $[a,s]=0$.
Zatem ciąg lokalizowany jest dokładny.
Ponieważ $K\otimes_R-$ jest naturalnie izomorficzny
z lokalizacją, $K$ jest modułem płaskim.
\end{proof}

\begin{przyklad}[Moduł płaski, który nie jest rzutowy]
\label{aa:prz:flat-not-projective}
Grupa addytywna $\mathbb Q$ jest płaskim
$\mathbb Z$-modułem na mocy poprzedniego twierdzenia.
Nie jest jednak modułem rzutowym. Gdyby nim była,
na mocy Twierdzenia~\ref{aa:tw:modules-projective-summand}
byłaby składnikiem sumy prostej pewnego modułu wolnego
$\mathbb Z^{(S)}$. W szczególności istniałby
niezerowy homomorfizm $\mathbb Q\to\mathbb Z^{(S)}$
będący włączeniem składnika.

Tymczasem każdy homomorfizm $h:\mathbb Q\to\mathbb Z$
jest zerowy: dla dowolnego $n\ge1$ mamy
$h(1)=n h(1/n)$, zatem liczba całkowita $h(1)$
jest podzielna przez każde $n$, co wymusza $h(1)=0$.
Dla $a\in\mathbb Z$, $b\ge1$ mamy dalej
$b\,h(a/b)=a\,h(1)=0$. Ponieważ $\mathbb Z$
nie ma torsji, $h(a/b)=0$ dla każdej liczby wymiernej.
Każda projekcja współrzędna homomorfizmu
$\mathbb Q\to\mathbb Z^{(S)}$ jest więc zerowa,
stąd sam homomorfizm jest zerowy. Otrzymaliśmy sprzeczność.
\end{przyklad}

\section{Powiązanie iloczynu tensorowego z modułami homomorfizmów}
\begin{twierdzenie}[Związek iloczynu tensorowego z homomorfizmami]
\label{aa:tw:tensor-hom-adjunction}
Niech $R$ będzie przemiennym pierścieniem z jedynką,
a $M,N,P$ modułami nad $R$. Istnieje naturalny
izomorfizm $R$-modułów
\[
\operatorname{Hom}_R(M\otimes_RN,P)\cong
\operatorname{Hom}_R\bigl(M,\operatorname{Hom}_R(N,P)\bigr).
\]
Homomorfizm $F:M\otimes_RN\to P$ odpowiada
odwzorowaniu $m\mapsto(n\mapsto F(m\otimes n))$.
\end{twierdzenie}
\begin{proof}
Niech $F:M\otimes_RN\to P$ będzie $R$-liniowy.
Dla ustalonego $m\in M$ odwzorowanie
$n\mapsto F(m\otimes n)$ jest $R$-liniowe
ze względu na dwuliniowość kanonicznego odwzorowania.
Wobec tego $\Phi(F)(m)$ należy do
$\operatorname{Hom}_R(N,P)$. Dla $m,m'\in M$
i $r\in R$ mamy
$\Phi(F)(m+m')=\Phi(F)(m)+\Phi(F)(m')$ oraz
$\Phi(F)(rm)=r\Phi(F)(m)$, co sprawdzamy na
każdym $n$. Zatem $\Phi(F)$ jest $R$-liniowe.

Odwrotnie, jeśli $G:M\to\operatorname{Hom}_R(N,P)$
jest $R$-liniowe, określmy $b(m,n)=G(m)(n)$.
Jest ono $R$-dwuliniowe: addytywność i zgodność
ze skalarami w pierwszej zmiennej wynika z
liniowości $G$, a w drugiej z liniowości $G(m)$.
Własność uniwersalna daje jedyny homomorfizm
$\Psi(G):M\otimes_RN\to P$ z
$\Psi(G)(m\otimes n)=G(m)(n)$.
Dla każdego $F$ i $G$ porównanie wartości na
wszystkich tensorach prostych pokazuje, że
$\Psi\Phi(F)=F$ oraz $\Phi\Psi(G)=G$.
Obie operacje zachowują dodawanie i skalary,
więc stanowią izomorfizm $R$-modułów.
\end{proof}

\begin{przyklad}[Obliczenie grupy homomorfizmów cyklicznych]
Dla dodatnich $m,n$ odwzorowanie
$f:\mathbb Z/m\mathbb Z\to\mathbb Z/n\mathbb Z$
jest wyznaczone przez $y=f([1])$, który musi spełniać
$my=0$. Niech $d=\gcd(m,n)$ i zapiszmy $m=dm'$, $n=dn'$,
gdzie $\gcd(m',n')=1$. Warunek $my=0$ oznacza,
że $n\mid m\widetilde y$, gdzie $\widetilde y\in\mathbb Z$
reprezentuje $y$. Jest on równoważny warunkowi
$n'\mid\widetilde y$. Zatem możliwe wartości $y$
są dokładnie klasami wielokrotności $n'$, tworzącymi
podgrupę cykliczną rzędu $d$. Dlatego
$\operatorname{Hom}_{\mathbb Z}(C_m,C_n)\cong C_d$.
Warto porównać ten rezultat z
Przykładem~\ref{aa:prz:tensor-cyclic}:
grupy otrzymane w tych dwóch obliczeniach
są izomorficzne, ale definicje konstrukcji są różne.
\end{przyklad}

\section{Związek z iloczynem tensorowym przestrzeni wektorowych}
\begin{stwierdzenie}[Baza iloczynu tensorowego skończenie wymiarowych przestrzeni]
\label{aa:st:tensor-vspace-dimension}
Niech $K$ będzie ciałem, a $V,W$ przestrzeniami wektorowymi
nad $K$ z bazami $(v_i)_{i=1}^m$ i $(w_j)_{j=1}^n$.
Wtedy rodzina $(v_i\otimes w_j)_{1\le i\le m,\,1\le j\le n}$
jest bazą $V\otimes_KW$; w szczególności
$\dim_K(V\otimes_KW)=mn$.
\end{stwierdzenie}
\begin{proof}
Każdy $v\in V$ i $w\in W$ ma postać
$v=\sum_i a_iv_i$, $w=\sum_j b_jw_j$.
Dwuliniowość daje
$v\otimes w=\sum_{i,j}a_ib_j(v_i\otimes w_j)$,
więc wskazana rodzina rozpina iloczyn tensorowy.

Dla niezależności niech $v_i^*:V\to K$ i
$w_j^*:W\to K$ oznaczają funkcjonały współrzędne,
zadane przez $v_i^*(v_k)=\delta_{ik}$ oraz
$w_j^*(w_\ell)=\delta_{j\ell}$.
Odwzorowanie $(v,w)\mapsto v_i^*(v)w_j^*(w)$
jest dwuliniowe. Z własności uniwersalnej wyznacza
funkcjonał $F_{ij}:V\otimes_KW\to K$, który
spełnia $F_{ij}(v_k\otimes w_\ell)=\delta_{ik}\delta_{j\ell}$.
Jeśli $\sum_{k,\ell}c_{k\ell}(v_k\otimes w_\ell)=0$,
to zastosowanie $F_{ij}$ daje $c_{ij}=0$.
Wszystkie współczynniki są zerowe, więc rodzina jest
liniowo niezależna i stanowi bazę.
\end{proof}

\begin{uwaga}[Co dalej?]
Iloczyn tensorowy jest prawostronnie dokładny, ale na ogół
nie zachowuje iniekcji. Dalszym algebraicznym narzędziem
badania tego zjawiska jest funktor $\operatorname{Tor}$,
zdefiniowany z wykorzystaniem rozdzielczości modułów.
Analogicznie funktor $\operatorname{Ext}$ bada dokładność
odwzorowań $\operatorname{Hom}$. Oba funktory zostaną
zdefiniowane i zbadane w części poświęconej algebrze
homologicznej; nie używamy ich tutaj bez uprzedniej definicji.
\end{uwaga}
